\documentclass[output=paper]{langscibook}
\ChapterDOI{10.5281/zenodo.22078567}
\title{Slavic semantics: A personal note}
\author{Hana Filip}
\abstract{}
\begin{document}
\maketitle
The start of my life-long preoccupation with semantics of Slavic languages goes back to 1989.  It is connected to two momentous events for me, which were surprisingly related, on December 29, 1989: namely, my first meeting with Barbara Partee at my first LSA (Linguistic Society of America) talk on the occasion of its annual meeting in Washington D.C., and Václav Havel's swearing-in ceremony as President of Czechoslovakia in Prague, the pinnacle of the Velvet Revolution that started on November 17 of that year. 

I was a third-year PhD student in the Department of Linguistics, at the University of California at Berkeley.  My talk, entitled \textit{Nominal reference and aspect: Aspectually determined sort-shifting in Czech}, was scheduled in the session Semantics II, at 3:50 pm. The analysis integrated insights from lexical semantics, in the style of my advisors Charles Fillmore and Paul Kay at UC Berkeley, with compositional semantics building on my previous studies at the Ludwig-Maximilians-University in Munich. The integration of
these two semantic traditions also characterized much of Barbara Partee's research, starting in the mid 1990s. Specifically, in her 
project on the Russian genitive of negation, funded by the 
National Science Foundation (NSF%\textsc{nsf}
) and inspired by her extended stays in Moscow, she integrated Russian (Moscow) lexical semantics with ``Western'' compositional semantics, particularly in the guise of Montague Semantics (she coined this term in her earliest introductions of Montague's work to linguists, 
see \citealt{Partee2005,Partee2015}).
Barbara Partee inspired many semanticists working on Slavic languages by showing how 
these two seemingly disparate semantic traditions can be successfully combined, leading 
to breakthroughs in our understanding of a variety of topics in Slavic semantics. 

Perhaps it was the mention of ``Czech'' in the title of my talk and the formal
semantics term ``sort-shifting'' that raised Barbara's attention. When I was getting ready to start my talk, she took the seat in the first row, right below the speaker's podium. That this raised the level of my jitters is an understatement. 
After my talk, she immediately came up to me: 
``May I introduce myself, I'm Barbara Partee.'' -- ``It's so nice to meet you, I know who you are.''  Barbara then started telling me with great enthusiasm and excitement that she had just returned from Prague, 
where she had spent the whole fall with Eva Hajičová, Petr Sgall, and their colleagues 
at the Institute of Formal and Applied Linguistics (Faculty of Mathematics and Physics, Charles University), as well as with philosophers Jaroslav Peregrin, Pavel Materna, and others, and  most importantly, she took part in the demonstrations during the Velvet Revolution in Prague. 
She directly experienced this seismic political and societal change in Czechoslovakia, and tremendously
enjoyed every moment of it. After my talk, and having just made acquaintance with Barbara Partee, I went back to my hotel room and switched on TV, the CNN news, showing President
Václav Havel walking across one of the courtyards at Prague Castle,
shortly after his swearing-in ceremony. It was one of the most memorable
moments of my life indeed, and one of the most iconic moments of the Czechoslovak history.
%(yes, and I did wonder, why are his pants so short, as many others did; in the excitement, chaos, and speed with which he got elected President, there was no time to fit his suit properly). 
For me it was also directly linked to the conversation that I just had with Barbara. 
This was the beginning of our friendship, Barbara mentored me 
and helped me along, like she did so many linguists. It has 
been her tireless, energetic and consistent support of many linguists who focus on Slavic semantics that was fundamental for this field to take off and flourish in the past almost forty years. 

% https://lsa.umich.edu/slavic/msp/fasl-proceedings-open-access.html

The initial decisive impetus came from the foundation of two conferences dedicated
to Slavic linguistics: namely, the \textit{Annual Workshops on Formal Approaches to Slavic Linguistics} (\textit{FASL%\textsc{fasl}
}) with its first meeting in 1992 at the University of Michigan, 
Ann Arbor, in the United States, and its European counterpart 
\textit{%The 
Formal Description of Slavic Languages} (\textit{FDSL%\textit{fdsl}
}) followed in 1995 with its first meeting at the University of Leipzig. 
The first FASL %\textsc{fasl} 
workshop was a rather small gathering of 
six Slavicists, John F. Bailyn, Steven Franks, Gerald Greenberg, Ljiljana Progovac, Maaike Schoorlemmer, and Jindřich Toman, 
who focused on \textit{Functional Categories in Slavic Syntax},
which is also the subtitle of the first proceedings, published in 1993 by Michigan Slavic Publications
%(Ann Arbor), 
and edited by Jindřich Toman (University of Michigan).
The second FASL %\textsc{fasl} 
workshop already had sixteen participants, and it was open to all the fields of Slavic linguistics, not just syntax. It took place at the Massachusetts
Institute of Technology (MIT%\textsc{mit}
) in 1993. The volume with papers based on the 
presented talks 
contained six semantics talks, five syntax talks, two morphology talks
(one delivered by Morris Halle) and one prosody talk. 
My paper in \textit{The MIT %\textsc{mit} 
Meeting} volume, entitled 
\textit{Quantificational Morphology: A Case Study in Czech}, was one
of my first publications on the topic of the generic \textit{-va-} in Czech; the most recent update on this topic is published in the present volume. The editors of \textit{The MIT %\textsc{mit} 
Meeting} volume, Sergey Avrutin (MIT%\textsc{mit}
), Steven Franks (Indiana University), and Ljiljana Progovac (Wayne State University), stated in their introduction: ``We believe that the rapid success of the \textit{Workshop on Formal Approaches to Slavic Linguistics} reflects an increased interest in this general area of linguistics. It is also our hope that the FASL %\textsc{fasl}
workshop will become a regular annual event, and that it will contribute to creating a sense of community and mutual awareness
among researchers from diverse backgrounds with similar professional interests.''
Their hope was fulfilled, because in 2025 the FASL %\textsc{fasl} 
workshop celebrated its thirty-fourth instantiation 
%in the Department of Linguistics 
at Cornell University, having enjoyed continued success throughout the years.
%, which attests to the high quality of the presented papers based on a highly competitive selection process. 

The inaugural conference of \textit{%The 
Formal Description of Slavic Languages} (\textit{FDSL}%\textsc{\textit{fdsl}}
) was organized in Leipzig in 1995 by Gerhild Zybatow, Uwe Junghanns, and Dorothee Fehrmann. Initially, it took place biennially, either at
the Leipzig University or the University of Potsdam, and later at other institutions, 
including the University of Nova Gorica (Slovenia), the Independent University of Moscow (Russia), Masaryk University in Brno (the Czech Republic), the University of Göttingen, Humboldt University of Berlin, the University of Graz, and the University of Wrocław. From 2006 to 2017, 
also ``halftime'' or \textit{x.5} conferences were organized at institutions outside Germany.
Starting in 2018, the conference became a regular annual event.
In 2025, FDSL18 was hosted by the University of Wrocław (Poland), where more than fifty presenters from across the world delivered papers spanning all the areas of Slavic linguistics. 
%on formal approaches to phonology, morphology, syntax and semantics, and many of which exploited the tools of corpus-based and experimental linguistics. Apart from the main session, the conference also had a workshop on \textit{Formal Approaches to Historical Slavic Linguistics}. The main session included contributions in the traditional areas of phonology, morphology, syntax and semantics.

Around the same time when the FASL %\textsc{fasl} 
workshop was initiated, one of its founders 
Steven Franks and his colleague George Fowler at the Indiana University
(Bloomington) launched a brand new journal, 
the \textit{Journal of Slavic Linguistics (JSL%\textsc{jsl}
)}. Their editorial manifesto in the first 1993 volume poses two ``unanswerable'' questions:
``Why are we devoting so much time and energy to \textit{JSL}%\textsc{jsl}
?  We have long felt that there is a real niche waiting to be filled by something like
\textit{JSL}%\textsc{jsl}
, and eventually came to ask ourselves `If not us, then who? And if
not now, then when?' Those questions are unanswerable, and so we
concluded that since there is no perfect occasion for launching a time and
energy-intensive endeavor, we might as well go ahead on our
own.''  Since 2006 the \textit{JSL} %\textsc{jsl} 
has also been the official journal of the \textit{Slavic Linguistics 
Society (SLS)}, which was established in December 2004,
%from a roundtable at the \textit{Annual Meeting of the American Association of Teachers of Slavic and East European Languages} (\textit{AATSEEL}) %\textsc{aatseel}), organized by Steven Franks, %and 
also thanks to Steven Franks, who served as its first chair, and 
thrives not the least thanks to its annual conference.
%, which promotes collaboration among linguists from a world-wide Slavic community.
For more than thirty years now, 
\textit{JSL} has been an important venue for disseminating state-of-the-art work on 
Slavic semantics. One of the \textit{JSL} Special Issues in 2003, guest edited by Wayles
Browne and Barbara Partee, is dedicated to Slavic semantics.  

This short note conveys some of my memories, and highlights only some
key milestones that early on fostered the research on Slavic semantics. This is not to say that many other key events and inspiring researchers did not contribute to it and without whom Slavic linguistics, and Slavic semantics, would not be what it is today. This note is by no means a comprehensive summary. It is my hope that this volume will 
be one piece in the rich mosaic of the work done in Slavic semantics, creating a sense of community among all researchers who share the excitement for its enigmas. \\

\noindent
Hana Filip, Düsseldorf, December 18, 2025. 

\printbibliography[heading=subbibliography,notkeyword=this]

\end{document}
